\documentclass[11pt]{article}
\usepackage[T1]{fontenc}\usepackage{lmodern}
\usepackage{amsmath,amssymb,amsthm,booktabs,hyperref,geometry,longtable}
\geometry{margin=2.3cm}
\newtheorem{theorem}{Theorem}
\title{A systematic filter for polynomial continued fractions with factorial reduction:\\ hidden Euler-type cases and fifteen irreducible candidates}
\author{Vichien Fugsukjit\\ Independent researcher, Bangkok, Thailand\\ \texttt{9chienvi9@gmail.com}}
\date{DRAFT, \today}
\begin{document}\maketitle
\begin{abstract}
We describe a reproducible pipeline that searches polynomial continued fractions (PCFs) with quadratic partial denominators and
quartic partial numerators, keeps those with full factorial reduction (FR), and removes every case explained by known mechanisms:
the family of Conjecture~1.3 of Ben David et al., deflation, rational or quadratic limits, M\"obius transforms of standard constants
and of the limits in H.~Cohen's databases, simple ${}_3F_2$ values, and Euler-type PCFs detected through hypergeometric solutions of
the three-term recurrence. A PCF that passed all numerical identification tests turns out to be of Euler type in disguise: its
numerators have the closed form $P_n=-(n+1)!\prod_{m=0}^{n-1}(3m+2)\,(n^2+3n+3)$, which gives its limit as an explicit series and
explains its gcd structure. Applying the full filter to 5,274 FR cases leaves fifteen PCFs whose associated order-four Fuchsian
operator is irreducible over $\mathbb Q(z)$ and whose limits match none of 21,984 tested constants at 1,000 digits. We publish them
as well-tested candidates, not as new constants, together with all code and data. Two further members of the same families have
limits that are M\"obius transforms of $\Gamma(1/3)$-type periods.
\end{abstract}

\section{Conventions}
A PCF is $x=a(0)+\cfrac{b(1)}{a(1)+\cfrac{b(2)}{a(2)+\cdots}}$ with $a,b\in\mathbb Z[n]$; $P_n,Q_n$ satisfy
$u_n=a(n)u_{n-1}+b(n)u_{n-2}$ with $P_{-1}=1$, $P_0=a(0)$, $Q_{-1}=0$, $Q_0=1$. The PCF has full FR if
$S_N=\frac1N\log(|Q_N|/\gcd(P_N,Q_N))$ stays bounded; numerically we require $|S_{800}-S_{200}|<0.15$ (exact integers).

\section{Search and filters}
\begin{enumerate}
\item Brute force over about $2.3\times10^9$ PCFs ($\deg a=2$, $b$ a product of four linear factors) with M\"obius hashing against
standard constants: 239 unidentified FR-candidates, 45 with full FR, 23 outside Conjecture~1.3 and a ${}_3F_2$ scan.
\item Extended search on the six $a(n)$ of these 23: $b(n)=s\prod_{i=1}^4(k_in+t_i)$, $1\le k_i\le6$, $|t_i|\le10$, leading
coefficient $B$ with $\alpha^2+4B$ a square ($\alpha$ the leading coefficient of $a$): 5,274 PCFs with full FR.
\item Removed: Conjecture~1.3 including index shifts $-2..2$ (846), slow convergence (15), rational or quadratic limits (313),
deflation (3,684); clustering by M\"obius equivalence left 347 new clusters; a Cohen-database limit (1), $\pi\sqrt3$ (1), ${}_3F_2$ (2).
\item Euler type: an order-one right factor of the recurrence over $\mathbb Q(n)$ (Petkov\v sek's algorithm via ore\_algebra):
61 of the extended-search cases and one of the original 23 (Section~3).
\item The operator $L=\theta^2(\theta-1)^2-z\,a(\theta+1)\theta^2-z^2\,b(\theta+2)$, $\theta=z\,d/dz$, annihilates
$\sum_n Q_nz^n/(n!)^2$; factoring over $\mathbb Q(z)$ with ore\_algebra, fifteen operators are irreducible, the others factor as
$[1,1,1,1]$, $[1,1,2]$ or $[1,3]$.
\end{enumerate}

\section{A hidden Euler-type example}
\begin{theorem}
For $a(n)=4n^2+10n-3$ and $b(n)=-n^2(n+4)(3n-4)$,
$P_n=-(n+1)!\,\prod_{m=0}^{n-1}(3m+2)\,(n^2+3n+3)$ and
\[x=\frac{1}{S-1/3},\qquad S=-\frac1{24}\sum_{k\ge1}\frac{k!\,(k+2)(k+3)(k+4)}{(2/3)_k\,3^k\,(k^2+k+1)(k^2+3k+3)} .\]
\end{theorem}
\begin{proof}[Proof sketch]
With $r(n)=P_n/P_{n-1}=(n+1)(3n-1)(n^2+3n+3)/(n^2+n+1)$ the identity $a(n)r(n-1)+b(n)=r(n)r(n-1)$ holds in $\mathbb Q(n)$, and the
initial values agree; the closed form follows by induction. The Casoratian of $(P,Q)$ gives
$Q_n/P_n=Q_0/P_0+\sum_{k=1}^{n}\prod_{j\le k}(-b(j))/(P_kP_{k-1})$, whose limit simplifies to $S-1/3$. Checked to 120 digits.
\end{proof}
The factor $\prod(3m+2)$ explains the gcd pattern (missing powers of $3$; the $p$-adic valuation depends on $3^{-1}\bmod p$), and
$n^2+3n+3=\tfrac14((2n+3)^2+3)$ is a norm form of $\mathbb Q(\sqrt{-3})$.

\section{Two \texorpdfstring{$\Gamma(1/3)$}{Gamma(1/3)}-type members}
With the same $a(n)=4n^2+10n-3$: for $b(n)=(n+3)(2n-5)(2n-1)(3n-1)$ the limit is a M\"obius transform of
$\Gamma(1/3)^3\sqrt3/\pi$ (match with a limit in Cohen's database); for $b(n)=(n+1)(2n-3)(2n+1)(3n-5)$ one has numerically
(250 digits) $x=(27-40c)/(9-24c)$ with $c=\pi^2/\Gamma(1/3)^3$. Both operators factor as $[1,3]$.

\section{The fifteen candidates}
{\small\begin{longtable}{lllll}\toprule
 & $a(n)$ & $b(n)$ & $x$ & $\rho$ \\\midrule
%%TABLE%%
\bottomrule\end{longtable}}
$\rho$: the non-integer local exponent of $L$ at the finite singularity of smallest modulus.

\paragraph{Tests passed by all fifteen (none identified).} Not in Conjecture~1.3 (shifts $-2..2$); not deflatable; limit not rational
or quadratic; not M\"obius ($|c|\le10^5$) of 36 standard constants; not M\"obius ($|c|\le6$) of the 635 limits of arXiv:2409.06086v4
or of the evaluable limit expressions of arXiv:2607.06581; not M\"obius ($|c|\le10^4$) of the 171 PCF limits of the euler2ai data
set; no hypergeometric solution over $\mathbb Q(n)$; operator irreducible over $\mathbb Q(z)$; PSLQ $[1,x,c,xc]$ ($|c|\le10^8$,
1,000 digits) against 21,984 constants (polylogarithms, Dirichlet $L$-values for $|d|\le100$, Clausen values, $\Gamma(j/d)$
products with $d\le24$, complete elliptic integrals at singular moduli); no match in the OEIS for the first 18 decimal digits of $|x|$ or for the first 8
reduced convergent denominators (controls: $\pi$ and Catalan's constant were found); the two $\Gamma(1/3)$ cases, used as positive controls, were
identified by the same code. For all fifteen, $x$ equals the ratio of the two connection coefficients of $L$ at its dominant
singularity; the individual coefficients are not $\Gamma$-products in an independent basis (150 digits; planted controls recovered).

\section{Limitations}
Novelty is not established: unpublished databases (e.g.\ the Ramanujan Library) were not available to us. The LMFDB has no search for
decimal constants and was not used. Irreducibility was tested over $\mathbb Q(z)$ only. A scheme-guided ${}_4F_3$ test failed its own control and is not counted. Full FR is a numerical criterion.

\section*{Data and code}
Scripts, logs and the candidate list (SHA-256 time-stamped in github.com/vichienFF/ramanujan-explorer before disclosure) are
released with this preprint.
\section*{Acknowledgements}
Computations used SageMath, ore\_algebra, ramanujantools, mpmath and Kaggle notebooks. The author used AI assistants (Anthropic
Claude for code and analysis, OpenAI GPT as an independent checker); all claims were verified by the released code.

\begin{thebibliography}{11}
\bibitem{RM} G.~Raayoni et al., Generating conjectures on fundamental constants with the Ramanujan Machine, Nature 590 (2021) 67--73.
\bibitem{PNAS} R.~Elimelech et al., Algorithm-assisted discovery of an intrinsic order among mathematical constants, PNAS 121 (2024), doi:10.1073/pnas.2321440121.
\bibitem{BD} N.~Ben David et al., On the connection between irrationality measures and polynomial continued fractions, Arnold Math.\ J.\ (2024), arXiv:2111.04468.
\bibitem{OD} O.~David, On Euler polynomial continued fractions, arXiv:2308.02567.
\bibitem{CMF} O.~David et al., The conservative matrix field, arXiv:2303.09318.
\bibitem{E2AI} T.~Raz et al., From Euler to AI: unifying formulas for mathematical constants, arXiv:2502.17533.
\bibitem{C1} H.~Cohen, A database of continued fractions of polynomial type, arXiv:2409.06086.
\bibitem{C2} H.~Cohen, Continued fractions of polynomial type: theory and encyclopedic dictionary, arXiv:2607.06581.
\bibitem{Hyper} M.~Petkov\v sek, Hypergeometric solutions of linear recurrences with polynomial coefficients, J.\ Symbolic Comput.\ 14 (1992) 243--264.
\bibitem{ore} M.~Kauers, M.~Jaroschek, F.~Johansson, Ore polynomials in Sage, LNCS 8942 (2015).
\bibitem{Apery} R.~Ap\'ery, Irrationalit\'e de $\zeta(2)$ et $\zeta(3)$, Ast\'erisque 61 (1979) 11--13.
\end{thebibliography}
\end{document}
